\documentclass[11pt]{article}

\usepackage[T1]{fontenc}
\usepackage{lmodern}
\usepackage[margin=1in]{geometry}
\usepackage{amsmath,amssymb,amsthm}
\usepackage{booktabs}
\usepackage{microtype}
\usepackage[hidelinks]{hyperref}
\hypersetup{
  pdftitle={The Least Uniquely Three-Colorable Integral Circulant with Clique Number Two},
  pdfauthor={Xiaojian Zeng},
  pdfsubject={Integral circulant graphs; unique three-colorability; exact finite census}
}
\emergencystretch=3em

\newtheorem{theorem}{Theorem}
\newtheorem{lemma}{Lemma}
\newcommand{\Z}{\mathbb Z}
\newcommand{\ICG}{\operatorname{ICG}}

\title{The Least Uniquely Three-Colorable Integral Circulant\\
with Clique Number Two}
\author{Xiaojian Zeng\\
\small PBC School of Finance, Tsinghua University, 100190\\
\small \texttt{xiaojianzeng@foxmail.com}}
\date{}

\begin{document}
\maketitle

\begin{abstract}
For a set $D$ of proper positive divisors of $n$, let $\ICG_n(D)$ be the
graph on $\Z_n$ in which distinct $x,y$ are adjacent when
$\gcd(x-y,n)\in D$.  We determine the least order in the restricted class
of uniquely $3$-colorable integral circulants having clique number $2$.
An exhaustive exact census gives order $60$, and at that order the only
divisor set is $D=\{2,5\}$.  Its unique unordered color partition is the
three residue classes modulo $3$.  The census covers all $4091$ pairs
$(n,D)$ with $n\leq 60$ and is reproduced by a separately implemented
CNF/DPLL verifier.  This result does not assert least order among uniquely
$k$-colorable integral circulants with arbitrary $k>\omega$.
\end{abstract}

\section{Statement and context}

A circulant graph is integral precisely when its connection set is a union
of gcd classes; this gives the divisor-set parametrization above
\cite{So2006}.  We say that a graph is \emph{uniquely $3$-colorable} when
its chromatic number is $3$ and all proper $3$-colorings induce the same
unordered partition into three nonempty independent sets.

Klotz and Sander asked whether every uniquely colorable circulant satisfies
$\chi=\omega$ \cite[Problem~3.6]{KlotzSander2017}.  Ba\v{s}i\'c answered that
question negatively by proving that
$\ICG_{12p}(\{2,p\})$ is uniquely $3$-colorable with clique number $2$ for
every prime $p\geq5$ \cite[Theorem~3.1]{Basic2026}.  In particular, $p=5$
supplies an order-$60$ example.  The issue addressed here is the sharp
finite boundary within the fixed parameter pair
$(\chi,\omega)=(3,2)$, together with the complete divisor-set
classification at that boundary.

\begin{theorem}\label{thm:main}
Let $n$ be a positive integer and let $D$ range over all subsets of the
proper positive divisors of $n$.  Among the graphs $\ICG_n(D)$ that are
uniquely $3$-colorable and have clique number $2$, the least possible order
is $n=60$.  At order $60$ the only such divisor set is
$D=\{2,5\}$.  Its unique unordered color partition is
\[
  C_r=\{x\in\Z_{60}:x\equiv r\pmod 3\},\qquad r=0,1,2.
\]
\end{theorem}

The minimum in Theorem~\ref{thm:main} is deliberately restricted to
uniquely three-colorable graphs with clique number two.  The theorem does
not exclude a smaller uniquely $k$-colorable integral circulant with
$k\geq4$ and $k>\omega$.

\section{Exact finite proof}

Write $\mathcal D_n$ for the set of proper positive divisors of $n$.
Every subset of $\mathcal D_n$ was visited, without an isomorphism or
symmetry quotient.  Thus the exact search domain has size
\[
 \sum_{n=1}^{59}2^{|\mathcal D_n|}=2043,
 \qquad
 \sum_{n=1}^{60}2^{|\mathcal D_n|}=4091.
\]
There is no duplication by equality of labeled graphs: if
$D\ne D'$ and $d\in D\mathbin\triangle D'$, then the edge $\{0,d\}$
belongs to exactly one of $\ICG_n(D)$ and $\ICG_n(D')$, since
$\gcd(d,n)=d$.

The certificate design follows a useful principle for finite-library
theorems: enumerate the entire discrete domain with exact predicates and
make the recorded witnesses independently reconstructible.  Related uses
of exhaustive graph certificates and a separate exact implementation occur,
in a different finite geometric setting, in \cite{ZengEtAl2026}.

For each pair $(n,D)$, the primary program constructs all adjacency rows as
integer bitsets directly from the gcd definition.  It then applies the
following exhaustive decision tree.

\begin{enumerate}
  \item An edgeless graph fails because its chromatic number is at most one.
  \item A disconnected graph fails unique $3$-colorability.
  \item A displayed triangle proves $\omega\geq3$.
  \item A displayed bipartition proves $\chi\leq2$.
  \item Every remaining graph is connected, has an edge, is triangle-free,
        and is nonbipartite.  Hence $\omega=2$ and $\chi\geq3$.
        A complete canonical DSATUR search then reports no proper
        $3$-coloring, two distinct unordered color partitions, or the
        unique unordered partition.
\end{enumerate}

For completeness, the second exclusion is justified even when a component
has chromatic number greater than three.  Such a graph is not
$3$-colorable.  Otherwise, if the whole graph is $3$-chromatic, at least one
component uses all three colors.  Permuting colors nontrivially on a
different nonempty component, while leaving the first component fixed,
changes the global unordered partition.  Thus a disconnected graph cannot
be uniquely $3$-colorable.

Using the saturation-degree ordering of Br{\'e}laz \cite{Brelaz1979}, the
search introduces colors only in their first-use order.  Its
vertex choice depends on saturation cardinality, degree, and vertex label,
not on color names.  Consequently every unordered proper $3$-color
partition appears in exactly one branch.  Stopping after two partitions is
sufficient only for the nonuniqueness decision; searches reporting zero or
one partition run to exhaustion.

The complete census is as follows.
\begin{center}
\begin{tabular}{lr}
\toprule
classification & number of pairs \\
\midrule
edgeless & 60 \\
disconnected & 677 \\
contains a triangle & 3228 \\
bipartite & 91 \\
not $3$-colorable & 17 \\
nonuniquely $3$-colorable & 17 \\
qualifying & 1 \\
\midrule
total & 4091 \\
\bottomrule
\end{tabular}
\end{center}
No pair with $n<60$ qualifies.  Among the $2048$ divisor sets at order
$60$, the sole qualifying record is $D=\{2,5\}$.  Its $360$ edges include
the odd cycle
\[
  4,2,0,14,9,4,
\]
and no triangle.  The exhaustive coloring certificate gives exactly the
three classes $C_0,C_1,C_2$ in Theorem~\ref{thm:main}.  These classes are
also visibly independent: a nonzero difference within one class is
divisible by $3$, so its gcd with $60$ is neither $2$ nor $5$.

An independent verifier reconstructs each graph using ordinary neighbor
sets, tests connectivity by union--find, detects triangles by set
intersection, and tests bipartiteness by parity propagation.  On every
remaining graph it forms the standard three-color CNF: each vertex has
exactly one color and adjacent vertices have different colors.  Vertex $0$
is fixed to color $0$.  After finding a model, the verifier blocks both
labeled models compatible with that choice that represent the same
unordered partition, and solves again by exhaustive unit propagation and
binary DPLL branching \cite{DavisLogemannLoveland1962}.  A second model is
precisely a second unordered
partition; unsatisfiability is precisely uniqueness.  This implementation
reproduced all $4091$ classifications, all witnesses, and the sole qualifier.
The domain count and the exhaustive alternatives above prove
Theorem~\ref{thm:main}.

\section{Reproducibility and assistance disclosure}

The ancillary files \texttt{evidence/enumerate\_icg.py} and
\texttt{evidence/verify\_census\_sat.py} use only the Python~3 standard
library.  From the project root, the census and its independent replay are
regenerated by
\begin{verbatim}
python3 evidence/enumerate_icg.py \
  --records evidence/census-primary.jsonl \
  --summary evidence/census-primary-summary.json
python3 evidence/verify_census_sat.py \
  --primary evidence/census-primary.jsonl \
  --records evidence/census-independent.jsonl \
  --audit evidence/census-audit.json
\end{verbatim}
The computations use exact integer and Boolean operations; no numerical
tolerance or floating-point inference enters the theorem.

An AI coding and reasoning assistant was used to help develop and inspect
the enumeration, verification code, and exposition.  The mathematical claim
rests on the explicitly described exhaustive algorithms and replayable exact
records; a separate implementation reconstructs every graph and repeats all
classification decisions.

\bibliographystyle{plain}
\bibliography{references}

\end{document}
